% Copyright (c) 2026 Eugene Joseph M. Ragasa.
% SPDX-License-Identifier: Apache-2.0
% Author: Eugene Joseph M. Ragasa
% Formal proof completed: 2026-10-19
\documentclass[11pt]{article}
\usepackage[T1]{fontenc}
\usepackage{amsmath,amssymb,booktabs,tabularx}
\usepackage[margin=1in]{geometry}
\title{Proof Dossier: Constant-Weight Rank-Two Frobenius Bridge}
\author{Eugene Joseph M. Ragasa}
\date{Formal proof completed: 2026-10-19}
\begin{document}\maketitle
\noindent\begin{tabularx}{\textwidth}{@{}lX@{}}\toprule
Statement ID & \texttt{BG-FOURIER-THM-029} \\
Status & proved; no admissions \\
Lean declaration & \texttt{ProjectKoios.BandGauge.Fourier.}\newline\texttt{rankTwoWeightedFrobeniusSq\_const\_eq\_coefficient\_mul\_power} \\
Owning module & \texttt{ProjectKoios/BandGauge/Fourier/RankTwoWeightedBridge.lean} \\
Axioms & \texttt{propext}, \texttt{Classical.choice}, \texttt{Quot.sound} \\
\bottomrule\end{tabularx}
\section{Statement and argument}
For a constant coordinate weight, the source-compatible rank-two weighted squared Frobenius quantity is exactly $\mathrm{weight}/(2E_{\mathrm{ref}}^2)$ times the unweighted finite matrix Frobenius power. Expand the product-index sum and factor the common coefficient.
\section{Boundary}
The constant-weight hypothesis is explicit and does not replace a nonuniform quadrature. Lean checked the exact statement; the dossier and proof were AI-assisted.
\end{document}
